\chapter{Research Methodology}
\label{ch:methodology}

\section{Research Approach}
The study is a retrospective, observational modelling study on de-identified
critical-care data, conducted as quantitative empirical research. The research
question (Section~\ref{sec:rq}) is operationalised as a prediction task with
pre-registered baselines (SIRS, qSOFA, SOFA) and an explicit null hypothesis for
the proportional-hazards assumption. A deliberate design choice constrains all
modelling to classical statistics: this preserves interpretability, supports
formal assumption testing, and establishes a transparent, transportable baseline
against which machine-learning approaches can later be compared on identical
data splits.

The methodology follows the standard prognostic-model development cycle
recommended by the TRIPOD (Transparent Reporting of a multivariable prediction
model for Individual Prognosis or Diagnosis) guidelines \cite{collins2015tripod}:
cohort definition,
outcome labelling, feature construction, model fitting, internal validation,
and external validation. Two overarching principles govern every stage.

\paragraph{Leakage safety.} At any prediction time, only information that
would genuinely be available to a clinician at that hour is used. This is
implemented through forward-filled values (each variable retains its most
recent measurement), time-since-last-measurement counters (encoding the
staleness of information as an explicit signal), and strict temporal
conditioning (the landmark design ensures that features are drawn only from
the observation window preceding the prediction point).

\paragraph{Honest evaluation.} Discrimination is measured out of sample using
patient-grouped cross-validation, ensuring that no patient contributes to both
the training and the test set. The frozen model is additionally transported to
a second, independent database with no refitting, providing the most stringent
form of external validation.

\section{Software Development Methodology}
The pipeline was developed incrementally as a sequence of numbered,
single-purpose stages (Chapter~\ref{ch:design}), each writing an inspectable
intermediate artefact that the next stage consumes. This staged design supports
reproducibility and debuggability: any stage can be re-run in isolation, and a
single configuration module is the sole source of truth for file paths,
clinical item identifiers, antibiotic name lists, and Sepsis-3 thresholds.
Data extraction and panel construction are implemented in Python (pandas);
statistical modelling and reporting are implemented in R, with each analysis
stage emitting a self-contained, reproducible report. All stochastic steps use
a fixed random seed (486649), and the multi-gigabyte source tables are processed
by streaming in bounded-memory chunks rather than being loaded whole.

\section{Data Sources}
\label{sec:data}

\subsection{Development Data: MIMIC-IV v3.1}
The primary development data is MIMIC-IV v3.1 \cite{johnson2023mimiciv}, a
large, de-identified critical-care database comprising records from the Beth
Israel Deaconess Medical Center (BIDMC) in Boston, Massachusetts. The database
contains comprehensive clinical data including demographics, vital signs,
laboratory measurements, medications, procedures, and clinical notes for over
65{,}000 adult ICU stays. It was accessed under the required PhysioNet
credentialing agreement.

\subsection{External Validation Data: eICU-CRD v2.0}
External validation uses the eICU Collaborative Research Database v2.0
\cite{pollard2018eicu}, a multi-centre collection spanning more than 200
United States ICUs. The eICU-CRD is independently collected, uses different
electronic health record systems, and represents a broader range of clinical
practices and patient populations than the single-centre MIMIC-IV. Both the
geographic and institutional diversity of eICU-CRD make it a meaningful
external test: if the model's signal is real, it should survive the transition
to this heterogeneous population.

\section{Cohort Construction and Inclusion Criteria}
An adult ICU cohort was constructed by applying the following inclusion
criteria:
\begin{itemize}
  \item Age $\ge 18$ years at ICU admission.
  \item First ICU stay only (to avoid within-patient correlation from repeated
  admissions).
  \item Length of stay $\ge 4$ hours (to ensure sufficient data for hourly
  feature construction).
\end{itemize}

This yielded \textbf{65{,}241 stays} in MIMIC-IV. The identical criteria were
applied to eICU-CRD, yielding a comparable external cohort.

\section{Outcome Definition: Sepsis-3 Onset}
\label{sec:outcome}

Sepsis-3 is defined following \textcite{seymour2016criteria}
as suspected infection accompanied by an acute rise of two or more points in
the SOFA score. Suspected infection is operationalised as a qualifying
antibiotic order paired with a body-fluid culture within the Sepsis-3 time
windows (culture within 24 hours before to 72 hours after the antibiotic, or
antibiotic within 72 hours of culture). Surveillance swabs (e.g., MRSA screens)
are excluded to avoid inflating the suspicion rate.

The cohort prevalence of Sepsis-3 on the first-24h analysis matrix is
\textbf{37.1\%}, with septic shock at 9.2\% and in-hospital mortality of 18.7\%
in the septic group versus 6.1\% in the non-septic group. These rates are
consistent with a valid label and with the published literature. A slightly
smaller fraction (32.0\%) can be assigned an \emph{exact hourly onset time}
within the 72-hour panel; this temporal subset is the one carried forward into
the landmark analyses, since it is the onset \emph{hour}, not merely the label,
that the real-time task requires.

\section{Hourly Panel and Temporal Alignment}
\label{sec:panel}

Every stay was resampled onto a regular hourly grid, yielding approximately
\textbf{3.1 million stay-hours} across the 65{,}241 stays (mean 47 hours per
stay). For each variable, three quantities are retained per hour:
\begin{enumerate}
  \item The \emph{raw value}: the measurement recorded at that hour, if any.
  \item The \emph{forward-filled value}: the most recent measurement, which is
  what the clinician currently knows.
  \item The \emph{time-since-last-measured} counter: the number of hours since
  the variable was last recorded, encoding irregular sampling as an explicit
  signal.
\end{enumerate}

Hourly SOFA is scored from the forward-filled values using the six organ
subsystems (respiration, coagulation, liver, cardiovascular, central nervous
system, and renal). The exact onset hour is located as the first hour inside
the suspected-infection window at which the hourly SOFA rise reaches the
two-point threshold. A per-hour binary label ``onset within 12 hours'' is
then derived for early-warning evaluation.

\section{The Landmark Design}
\label{sec:landmark}

The \textbf{hour-six landmark} is central to the study design. At hour 6 of
each stay, the analysis conditions on the patient being still at risk
(not yet septic), yielding \textbf{48{,}829 at-risk stays}. Of these,
\textbf{4{,}750} (9.7\%) develop incident sepsis within the subsequent
observation window, and \textbf{3{,}330} (6.8\%) within 12 hours. This is a
realistically imbalanced early-prediction problem, unlike the near-balanced
classification that a static cross-sectional analysis produces. It also
separates the clinically actionable target (incident sepsis) from prevalent
sepsis that is already present at admission and not predictable by any
early-warning system.

For the dynamic landmark supermodel (Analysis~2), the landmark construction is
repeated at hours $\{6, 12, 18, 24, 36, 48\}$, producing a stacked dataset of
\textbf{207{,}010 stay-landmark rows} across 48{,}829 unique stays. The
supermodel is fitted once on this stacked dataset with the landmark time $s$,
its quadratic term $s^2$, and all feature-by-$s$ interactions as additional
covariates, so that evaluating the model at the current hour yields
continuously updated risk.

\section{Statistical Models}
\label{sec:models}

The pipeline implements seven analytical stages, each addressing a specific
methodological question. The first analysis establishes the discrimination
ceiling; the second builds the real-time deliverable; analyses three through
seven provide complementary inferential and evaluative evidence.

\paragraph{Analysis 1: Full-feature elastic-net nomogram.}
An elastic-net penalised logistic regression at the hour-six landmark uses the
full feature set: vitals (heart rate, respiratory rate, MAP, SpO$_2$,
temperature, GCS), severity scores (SOFA, qSOFA, SIRS), laboratory values
(lactate, creatinine, WBC, P/F ratio), and six-hour trajectory slopes. Missing
labs (70--88\% missing at hour 6) are handled by an informative-missingness
design: median imputation plus a binary ``was measured'' indicator per lab and
per slope. Evaluation uses nested cross-validation (5-fold outer, 5-fold inner
for $\lambda$ selection).

\paragraph{Analysis 2: Dynamic landmark supermodel (central contribution).}
The stacked landmark dataset is fitted as one penalised logistic model with
$s$, $s^2$, and 27 feature-by-$s$ interactions, yielding 56 candidate features.
Patient-grouped folds prevent information leakage across landmarks. The
single model emits hourly-updating risk at any prediction hour, a genuine
real-time deliverable built entirely from classical statistics.

\paragraph{Analysis 3: Within-patient correlation.}
Generalised estimating equations (GEE) with an exchangeable working correlation
structure and cluster-robust (sandwich) standard errors quantify the
within-patient correlation ignored by pooled models. A ridge-penalised
time-varying Cox model on the counting-process formulation provides
inferentially valid hazard ratios.

\paragraph{Analysis 4: Real-time evaluation.}
The supermodel is evaluated on a held-out set of 6{,}000 stays using the
PhysioNet/CinC 2019 utility score and an alarm-burden analysis that trades
sensitivity against false alarms per patient-day.

\paragraph{Analysis 5: Subgroup fairness.}
Performance is stratified by sex and age band along three axes: discrimination
(AUROC), calibration-in-the-large, and error-rate parity (sensitivity and
false-positive rate at a shared decision threshold), the last giving an
equal-opportunity view of equitable performance across demographic subgroups.

\paragraph{Analysis 6: Group-based trajectory modelling.}
A latent-class growth model is fitted to first-24h SOFA trajectories, and
class membership is linked to both mortality and incident sepsis via Cox
regression and logistic regression respectively.

\paragraph{Analysis 7: Confounder-adjusted causal analysis.}
The effect of early antibiotic administration ($\le 3$ hours) on in-hospital
mortality among Sepsis-3 patients is estimated using four estimators of
increasing robustness: crude unadjusted, propensity-score matching (PSM),
inverse-probability weighting (IPW), and augmented IPW (doubly robust).

\section{Evaluation Framework}
\label{sec:eval_framework}

\paragraph{Discrimination.} The area under the receiver operating
characteristic curve (AUROC) with DeLong 95\% confidence intervals.
Precision-recall AUC (PR-AUC) supplements AUROC for the imbalanced
incident-onset task.

\paragraph{Calibration.} The Brier score, calibration-in-the-large (mean
predicted risk versus observed event rate), and the calibration slope (from a
logistic recalibration model; slope = 1 indicates perfect calibration of the
linear predictor's spread).

\paragraph{Clinical utility.} The PhysioNet/CinC 2019 normalised utility
score (0 = never-alarm baseline, 1 = per-hour oracle) and an alarm-burden
curve plotting timely detection sensitivity against false alarms per
patient-day across a range of risk thresholds.

\paragraph{External validation.} The frozen MIMIC coefficients are applied
without any refitting to the eICU-CRD hour-six landmark table. The
frozen-transport AUROC is compared against an eICU-native elastic-net refit
(the internal eICU ceiling) to separate loss of discrimination from
miscalibration.
